\section{Physical $n=2$ metric KKT projection}
\label{sec:bhsm-metric-kkt}

The remaining cosmological action object can be written explicitly in a
two-variable scalar metric basis once the momentum constraint and scalar
longitudinal gauge have been handled.  Later BHSM fixed-boundary work supplies
an exact action-derived domain architecture for this reduction, although its
radial coefficients belong to a different background and are not imported
numerically into cosmology.

\subsection{Action-derived domain architecture}

For a nonzero scalar harmonic, the induced metric trace before scalar gauge
fixing has two components,
\begin{equation}
\Theta
=
\psi(1)+\frac14\Box_4E(1),
\qquad
\mathcal E=E(1).
\end{equation}
The radial lapse $A$ and scalar normal--tangential shift $B$ are not induced
metric trace data.  The action-derived fixed-$h$ construction proceeds in the
following order:

\begin{enumerate}
\item vary before gauge fixing;
\item use the momentum equation to eliminate the shift sector;
\item choose the scalar longitudinal gauge $E=0$;
\item impose the induced-metric Dirichlet trace
      $\psi(1)=0$;
\item retain the matcher reaction
      $\eta_{\rm tr}=P_\psi(1)$;
\item quotient endpoint-preserving radial diffeomorphisms.
\end{enumerate}

The extended Green current contains the finite matcher pairing
\begin{equation}
\psi_1\eta_2-\psi_2\eta_1,
\end{equation}
which cancels the canonical-momentum endpoint current on the constrained
domain.  This gives equal operator and adjoint domains in the corresponding
indefinite KKT saddle problem.

\paragraph{Domain proposition.}
The cosmological $n=2$ metric calculation adopts this sequence as the
\emph{action-derived domain architecture candidate}
$\mathcal D_{g,\rm phys}^{(2)}$.  The domain architecture is inherited; the
background-dependent radial coefficients are not.

\subsection{Two-variable physical metric block}

After the above reduction, let
\begin{equation}
y_g
=
\begin{pmatrix}
A\\
\psi
\end{pmatrix}
\end{equation}
denote the scalar metric coordinates.  The physical metric Hessian projected
onto the $n=2$ harmonic is
\begin{equation}
\boxed{
\mathcal A_g^{(2)}
=
\begin{pmatrix}
a_{AA} & a_{A\psi}\\
a_{A\psi} & a_{\psi\psi}
\end{pmatrix},
}
\label{eq:Ag2-explicit}
\end{equation}
where each entry is an action matrix element evaluated on
$\mathcal D_{g,\rm phys}^{(2)}$.

The critical-fold principal operator previously derived in BHSM demonstrates
that a lapse--Weyl saddle structure of this type is action generated, but its
numerical coefficients are branch specific.  They are not substituted into
Eq.~\eqref{eq:Ag2-explicit}.

\subsection{Projected metric--topographic source}

Let
\begin{equation}
X_2
=
\begin{pmatrix}
q_2\\
\Pi_2
\end{pmatrix}.
\end{equation}
The mixed physical Hessian defines the source map
\begin{equation}
\boxed{
\mathcal B_{g2}
=
\begin{pmatrix}
b_{Aq} & b_{A\Pi}\\
b_{\psi q} & b_{\psi\Pi}
\end{pmatrix}.
}
\label{eq:Bg2-explicit}
\end{equation}

The retained BHSM scalar/topographic action already supplies a nonzero
metric--topographic stress-mixing second variation.  Projecting that bilinear
form onto the physical $n=2$ mode and the two metric basis directions defines
the entries in Eq.~\eqref{eq:Bg2-explicit}.  In particular, for metric test
variation $\gamma$ and topographic test variation $\phi$, the source contains
terms of the form
\begin{equation}
B_{g\phi}[\gamma,\phi]
=
\int d\mu_g\,w\,\gamma^{\mu\nu}
\left[
2F''(X)(j\!\cdot\!d\phi)A_{\mu\nu}
+
2F'(X)j_{(\mu}\partial_{\nu)}\phi
-
F'(X)g_{\mu\nu}(j\!\cdot\!d\phi)
\right].
\label{eq:metric-topographic-mixed-form}
\end{equation}
Thus the remaining coefficients are matrix elements of a stored action
variation, not independent phenomenological couplings.

\subsection{Explicit response formula}

Define
\begin{equation}
\Delta_g
=
a_{AA}a_{\psi\psi}-a_{A\psi}^2.
\end{equation}
If $\Delta_g\neq0$, then
\begin{equation}
(\mathcal A_g^{(2)})^{-1}
=
\frac1{\Delta_g}
\begin{pmatrix}
a_{\psi\psi} & -a_{A\psi}\\
-a_{A\psi} & a_{AA}
\end{pmatrix},
\end{equation}
and the physical metric response is
\begin{equation}
\boxed{
\mathcal R_{g2}
=
-(\mathcal A_g^{(2)})^{-1}\mathcal B_{g2}.
}
\label{eq:Rg2-explicit}
\end{equation}

Therefore all metric response amplitudes are fixed once the seven independent
action matrix elements
\begin{equation}
\boxed{
a_{AA},\;
a_{A\psi},\;
a_{\psi\psi},\;
b_{Aq},\;
b_{A\Pi},\;
b_{\psi q},\;
b_{\psi\Pi}
}
\label{eq:seven-matrix-elements}
\end{equation}
are evaluated on the physical domain.

\subsection{Schur consistency with the scalar/topographic block}

With the candidate scalar/topographic block
$\mathcal C_{2,\rm ST}$ from
Sec.~\ref{sec:closed-flrw-metric-tangent}, the reduced mode operator is
\begin{equation}
\boxed{
\mathcal K_2
=
\mathcal C_{2,\rm ST}
-
\mathcal B_{g2}^{\dagger}
(\mathcal A_g^{(2)})^{-1}
\mathcal B_{g2}.
}
\label{eq:K2-explicit-metric}
\end{equation}
Equations~\eqref{eq:Rg2-explicit} and
\eqref{eq:K2-explicit-metric} show that the observable metric response and
the mode self-energy correction are controlled by the same finite set of
action matrix elements.

\subsection{Updated closure count}

The action-provenance problem is now reduced to two tasks:

\begin{enumerate}
\item confirm that the fixed-$h$ matcher/Green-current architecture above is
      the correct physical domain for the closed-FLRW $n=2$ branch;
\item evaluate the seven matrix elements in
      Eq.~\eqref{eq:seven-matrix-elements} from the retained action.
\end{enumerate}

No additional observable-specific anisotropic coefficient is permitted.
